\documentclass[12pt]{article}
\usepackage[utf8]{inputenc}
\usepackage[T1]{fontenc}
\usepackage[margin=1in]{geometry}
\usepackage{setspace}
\usepackage{xcolor}
\usepackage[colorlinks=true, urlcolor=blue!60!black, linkcolor=blue!60!black]{hyperref}
\setlength{\parindent}{0em}
\setlength{\parskip}{0.6em}
\renewcommand{\baselinestretch}{1.0}
\begin{document}
\begin{center}
{\LARGE \textit{Nevada Department of Human Resources v. Hibbs}, 538 U.S. 721 (2003)\par}
\vspace{0.5em}
{\large Case Brief --- Two-Page Edition\par}
\end{center}

\section{Memory Jogger}
Congress validly abrogated state sovereign immunity for the FMLA's family-care leave provision, because it targets sex-role stereotyping subject to heightened scrutiny with a congruent and proportional remedy (\href{https://supreme.justia.com/cases/federal/us/538/721/}{Justia}).

\section{Detailed Case Facts}
The FMLA grants 12 weeks of unpaid leave to care for a seriously ill family member and authorizes damages suits against public employers. Nevada fired William Hibbs after he exhausted 12 weeks of leave taken to care for his injured wife and did not return.

\section{Procedural History}
The District of Nevada held the suit barred by the Eleventh Amendment; the Ninth Circuit reversed, splitting with the Fifth Circuit.

\section{Judicial Votes}
Affirmed, 6--3. Rehnquist for the Court (O'Connor, Souter, Ginsburg, Breyer); Souter concurring; Stevens concurring in the judgment; Scalia and Kennedy (with Scalia and Thomas) dissenting.

\section{Holding}
Under \textit{Boerne}'s congruence-and-proportionality test, the family-care provision validly abrogates immunity. Because sex classifications get heightened scrutiny, Congress can more easily show a pattern of unconstitutional state conduct than for age or disability. The self-care provision was not addressed (later barred in \textit{Coleman}).

\section{Analysis of Opinions}
\textbf{Rehnquist:} History and the FMLA record showed state leave policies built on stereotypes that women are caregivers and men are not, with 15 States offering women extended maternity leave and only four offering men parallel leave. A gender-neutral floor, with unpaid leave, eligibility limits, and capped damages, was congruent and proportional. \textbf{Stevens:} would rest on Congress's Commerce power. \textbf{Scalia:} ``[T]here is no guilt by association'' among States. \textbf{Kennedy:} the record was thin and the FMLA is a benefit program, not a remedy.

\section{Examples: Future Applications}
\textbf{Same-Side:} ADA Title II suits over courthouse access (\textit{Tennessee v. Lane}). \textbf{Opposite-Side:} FMLA self-care suits against States (\textit{Coleman}). \textbf{Fence-Sitter:} Damages suits under a pregnancy-accommodation statute, given \textit{Geduldig}.

\section{Critique}
Federalism critics say \textit{Hibbs} relaxed \textit{Boerne} on thin evidence. Progressive scholars praise its ``feminist universalism'' but fault a scrutiny-based hierarchy that leaves older and disabled workers unprotected.

\section{Key Quotations}
\begin{quote}
``Stereotypes about women's domestic roles are reinforced by parallel stereotypes presuming a lack of domestic responsibilities for men. Because employers continued to regard the family as the woman's domain, they often denied men similar accommodations or discouraged them from taking leave. These mutually reinforcing stereotypes created a self-fulfilling cycle of discrimination that forced women to continue to assume the role of primary family caregiver, and fostered employers' stereotypical views about women's commitment to work and their value as employees.''
\end{quote}
\noindent --- Chief Justice Rehnquist, Opinion of the Court, 538 U.S. at 736.

\end{document}
